
\documentclass[aspectratio=169]{beamer}

\usetheme{Madrid}
\usecolortheme{default}

\usepackage{amsmath,amssymb,amsfonts,bm,mathtools}
\usepackage{physics}
\usepackage{listings}
\usepackage{xcolor}

\definecolor{codegray}{rgb}{0.95,0.95,0.95}
\lstset{
    basicstyle=\ttfamily\tiny,
    backgroundcolor=\color{codegray},
    breaklines=true,
    frame=single,
    columns=fullflexible,
    showstringspaces=false,
    keywordstyle=\color{blue},
    commentstyle=\color{teal},
    stringstyle=\color{red}
}

\title{Loss Functions for Recurrent Neural Networks on Spatiotemporal Grids}
\subtitle{Including an Example of a Physics-Informed RNN in PyTorch}
\author{Morten Hjorth-Jensen}
\date{FYS5429/9429}

\begin{document}

\frame{\titlepage}

\begin{frame}{Problem setting}
We consider a recurrent neural network for temporal data on a spatial grid
\[
(x,y)\in \Omega_x\times\Omega_y
\]
with two height fields (or layers)
\[
h_t^{(1)}(x,y), \qquad h_t^{(2)}(x,y).
\]

At each time \(t\), the target field is
\[
Y_t(x,y)=\bigl(h_t^{(1)}(x,y),\,h_t^{(2)}(x,y)\bigr).
\]

The central question is:
\begin{quote}
What is the most reasonable loss function for training such a recurrent model?
\end{quote}
\end{frame}

\begin{frame}{Main recommendation}
A good default choice is to use \textbf{one joint loss} over all times, grid points, and both layers,
rather than two completely separate training objectives.

\medskip

The reason is that the two layers are usually:
\begin{itemize}
\item measured on the same spatial grid,
\item coupled through the underlying dynamics,
\item and likely to share latent spatiotemporal structure.
\end{itemize}

Thus the model should usually output two channels and be trained jointly.
\end{frame}

\begin{frame}{A natural baseline loss}
A standard baseline is a weighted mean-squared error:
\[
\mathcal{L}
=
\frac{1}{T N_x N_y}
\sum_{t=1}^{T}\sum_{x,y}
\left[
w_1\bigl(\hat h_t^{(1)}(x,y)-h_t^{(1)}(x,y)\bigr)^2
+
w_2\bigl(\hat h_t^{(2)}(x,y)-h_t^{(2)}(x,y)\bigr)^2
\right].
\]

Equivalently,
\[
\mathcal{L} = w_1 \mathcal{L}_1 + w_2 \mathcal{L}_2,
\]
where \(\mathcal{L}_1\) and \(\mathcal{L}_2\) are the losses for the two layers.

This is still best viewed as a \emph{single joint objective}.
\end{frame}

\begin{frame}{Should one use two separate losses?}
Using two separate losses is usually not necessary unless:
\begin{itemize}
\item the two layers have very different scales,
\item one layer is more important scientifically,
\item one layer is much noisier,
\item or one wants to formulate a multi-task learning problem explicitly.
\end{itemize}

In practice, it is usually sufficient to define
\[
\mathcal{L} = \alpha \mathcal{L}_1 + \beta \mathcal{L}_2
\]
and tune \(\alpha,\beta\) or normalize the two channels beforehand.
\end{frame}

\begin{frame}{Why channel weighting matters}
If the magnitudes of the two layers differ strongly, then an unweighted MSE may be dominated by the larger field.

A practical choice is
\[
\alpha = \frac{1}{\sigma_1^2},
\qquad
\beta = \frac{1}{\sigma_2^2},
\]
where \(\sigma_1,\sigma_2\) are characteristic scales or standard deviations of the two layers.

This leads to a normalized objective in which both channels influence training more comparably.
\end{frame}

\begin{frame}{Alternative pointwise losses}
Besides MSE, one may use:
\begin{itemize}
\item \textbf{MAE:}
\[
\mathcal{L}_{\mathrm{MAE}} = \frac{1}{N}\sum |\hat y-y|
\]
which is more robust to outliers,
\item \textbf{Huber loss,}
which interpolates between MSE and MAE,
\item weighted losses if uncertainty or heteroscedastic noise is known.
\end{itemize}

For smooth physical fields, MSE is usually the best starting point.
\end{frame}

\begin{frame}{Why pointwise loss may not be enough}
For spatiotemporal fields, a pointwise MSE can produce overly smooth predictions.

If preserving coherent structures matters, it is often useful to add a spatial derivative penalty:
\[
\mathcal{L}
=
\mathcal{L}_{\mathrm{data}}
+
\lambda_{\nabla}\mathcal{L}_{\nabla},
\]
with
\[
\mathcal{L}_{\nabla}
=
\sum_{t,x,y,c}
\abs{\nabla \hat h_t^{(c)}(x,y)-\nabla h_t^{(c)}(x,y)}^2.
\]

This encourages the correct spatial gradients and sharp structures.
\end{frame}

\begin{frame}{Temporal weighting}
If the recurrent network predicts many future time steps, it is natural to use
\[
\mathcal{L}
=
\sum_{t=1}^{T}\omega_t \mathcal{L}_t,
\]
where \(\mathcal{L}_t\) is the loss at time \(t\).

Choices include:
\begin{itemize}
\item \(\omega_t=1\): equal weighting,
\item larger \(\omega_t\) for later times if long-term prediction is important,
\item larger \(\omega_t\) for early times if short-term stabilization is the main goal.
\end{itemize}
\end{frame}

\begin{frame}{A useful practical combined objective}
A practical recommendation is
\[
\mathcal{L}
=
\alpha \mathcal{L}_1
+
\beta \mathcal{L}_2
+
\lambda_{\nabla}\mathcal{L}_{\nabla}
+
\lambda_{\mathrm{phys}}\mathcal{L}_{\mathrm{phys}},
\]
where:
\begin{itemize}
\item \(\mathcal{L}_1,\mathcal{L}_2\): data losses for the two layers,
\item \(\mathcal{L}_{\nabla}\): gradient-matching loss,
\item \(\mathcal{L}_{\mathrm{phys}}\): physics-informed penalty.
\end{itemize}

This is often better than training the two layers independently.
\end{frame}

\begin{frame}{Physics-informed RNN idea}
A physics-informed RNN augments the ordinary data-fitting objective with a residual enforcing the governing equations.

Suppose the fields satisfy symbolic PDEs
\[
\partial_t h^{(1)} = \mathcal{F}_1(h^{(1)},h^{(2)}),
\qquad
\partial_t h^{(2)} = \mathcal{F}_2(h^{(1)},h^{(2)}).
\]

Then we define residuals
\[
R_1 = \partial_t \hat h^{(1)} - \mathcal{F}_1(\hat h^{(1)},\hat h^{(2)}),
\qquad
R_2 = \partial_t \hat h^{(2)} - \mathcal{F}_2(\hat h^{(1)},\hat h^{(2)}),
\]
and penalize them in the loss.
\end{frame}

\begin{frame}{Physics-informed loss}
A general physics-informed objective can be written as
\[
\mathcal{L}_{\mathrm{total}}
=
\mathcal{L}_{\mathrm{data}}
+
\lambda_{\mathrm{phys}}\mathcal{L}_{\mathrm{phys}},
\]
where
\[
\mathcal{L}_{\mathrm{data}}
=
\alpha \mathcal{L}_1 + \beta \mathcal{L}_2,
\]
and
\[
\mathcal{L}_{\mathrm{phys}}
=
\frac{1}{T N_x N_y}
\sum_{t,x,y}
\left(
R_1(t,x,y)^2 + R_2(t,x,y)^2
\right).
\]

This encourages both agreement with the data and consistency with the underlying PDE.
\end{frame}

\begin{frame}{Discrete derivatives on a grid}
On a uniform grid one may approximate derivatives by finite differences, for example
\[
\partial_t h_t \approx \frac{h_{t+1}-h_t}{\Delta t},
\]
\[
\partial_x h_{i,j} \approx \frac{h_{i+1,j}-h_{i-1,j}}{2\Delta x},
\qquad
\partial_y h_{i,j} \approx \frac{h_{i,j+1}-h_{i,j-1}}{2\Delta y}.
\]

These discrete operators can be implemented directly in PyTorch using tensor slicing.
\end{frame}

\begin{frame}{Recommended training strategy}
A good practical training strategy is:
\begin{enumerate}
\item normalize each layer separately,
\item use one recurrent model with two output channels,
\item start with weighted MSE over both channels,
\item add a physics term only after the baseline trains stably,
\item optionally add spatial-gradient penalties if the fields become too smooth.
\end{enumerate}
\end{frame}

\begin{frame}{Recommended default choice}
A strong default objective is
\[
\mathcal{L}
=
\frac{1}{T N_x N_y}
\sum_{t,x,y}
\sum_{c=1}^{2}
w_c\bigl(\hat h_t^{(c)}(x,y)-h_t^{(c)}(x,y)\bigr)^2
+
\lambda_{\mathrm{phys}}\mathcal{L}_{\mathrm{phys}}.
\]

So the answer is:
\begin{itemize}
\item usually \textbf{one joint loss},
\item but with \textbf{channel weights} if needed,
\item and possibly with a \textbf{physics-informed residual} term.
\end{itemize}
\end{frame}

\begin{frame}[fragile]{PyTorch example: physics-informed RNN (1)}
\begin{lstlisting}[language=Python]
import torch
import torch.nn as nn
import torch.nn.functional as F

class PhysicsInformedRNN(nn.Module):
    def __init__(self, nx, ny, hidden_size=128):
        super().__init__()
        self.nx = nx
        self.ny = ny
        self.input_size = 2 * nx * ny
        self.hidden_size = hidden_size
        self.rnn = nn.GRU(
            input_size=self.input_size,
            hidden_size=hidden_size,
            batch_first=True
        )
        self.out = nn.Linear(hidden_size, self.input_size)

    def forward(self, x):
        # x: [batch, time, 2, nx, ny]
        b, t, c, nx, ny = x.shape
        x_flat = x.reshape(b, t, c * nx * ny)
        h, _ = self.rnn(x_flat)
        y = self.out(h)
        y = y.reshape(b, t, 2, nx, ny)
        return y
\end{lstlisting}
\end{frame}

\begin{frame}[fragile]{PyTorch example: finite-difference operators (2)}
\begin{lstlisting}[language=Python]
def dx(field, dx_val):
    # field: [batch, time, nx, ny]
    return (field[:, :, 2:, 1:-1] - field[:, :, :-2, 1:-1]) / (2.0 * dx_val)

def dy(field, dy_val):
    return (field[:, :, 1:-1, 2:] - field[:, :, 1:-1, :-2]) / (2.0 * dy_val)

def dt(field, dt_val):
    # central time difference for interior points
    return (field[:, 2:, :, :] - field[:, :-2, :, :]) / (2.0 * dt_val)
\end{lstlisting}
\end{frame}

\begin{frame}[fragile]{PyTorch example: a simple PDE residual (3)}
\begin{lstlisting}[language=Python]
def physics_residual(pred, dt_val, dx_val, dy_val, nu=0.01):
    # Example residual for two coupled diffusion-like fields
    # pred: [batch, time, 2, nx, ny]
    h1 = pred[:, :, 0]
    h2 = pred[:, :, 1]

    h1_t = dt(h1, dt_val)
    h2_t = dt(h2, dt_val)

    h1_xx = (h1[:, 1:-1, 2:, 1:-1] - 2*h1[:, 1:-1, 1:-1, 1:-1]
             + h1[:, 1:-1, :-2, 1:-1]) / (dx_val**2)
    h1_yy = (h1[:, 1:-1, 1:-1, 2:] - 2*h1[:, 1:-1, 1:-1, 1:-1]
             + h1[:, 1:-1, 1:-1, :-2]) / (dy_val**2)

    h2_xx = (h2[:, 1:-1, 2:, 1:-1] - 2*h2[:, 1:-1, 1:-1, 1:-1]
             + h2[:, 1:-1, :-2, 1:-1]) / (dx_val**2)
    h2_yy = (h2[:, 1:-1, 1:-1, 2:] - 2*h2[:, 1:-1, 1:-1, 1:-1]
             + h2[:, 1:-1, 1:-1, :-2]) / (dy_val**2)

    r1 = h1_t[:, :, 1:-1, 1:-1] - nu * (h1_xx + h1_yy)
    r2 = h2_t[:, :, 1:-1, 1:-1] - nu * (h2_xx + h2_yy)
    return r1, r2
\end{lstlisting}
\end{frame}

\begin{frame}[fragile]{PyTorch example: total loss (4)}
\begin{lstlisting}[language=Python]
def total_loss(pred, target, dt_val, dx_val, dy_val,
               w1=1.0, w2=1.0, lambda_phys=1.0):
    # pred, target: [batch, time, 2, nx, ny]
    mse1 = F.mse_loss(pred[:, :, 0], target[:, :, 0])
    mse2 = F.mse_loss(pred[:, :, 1], target[:, :, 1])
    data_loss = w1 * mse1 + w2 * mse2

    r1, r2 = physics_residual(pred, dt_val, dx_val, dy_val)
    phys_loss = r1.pow(2).mean() + r2.pow(2).mean()

    return data_loss + lambda_phys * phys_loss, {
        "data_loss": data_loss.detach(),
        "phys_loss": phys_loss.detach(),
        "mse1": mse1.detach(),
        "mse2": mse2.detach(),
    }
\end{lstlisting}
\end{frame}

\begin{frame}[fragile]{PyTorch example: training step (5)}
\begin{lstlisting}[language=Python]
model = PhysicsInformedRNN(nx=32, ny=32, hidden_size=128)
optimizer = torch.optim.Adam(model.parameters(), lr=1e-3)

for batch_x, batch_y in train_loader:
    optimizer.zero_grad()

    pred = model(batch_x)
    loss, info = total_loss(
        pred, batch_y,
        dt_val=0.1,
        dx_val=1.0,
        dy_val=1.0,
        w1=1.0,
        w2=1.0,
        lambda_phys=0.1
    )

    loss.backward()
    torch.nn.utils.clip_grad_norm_(model.parameters(), 1.0)
    optimizer.step()
\end{lstlisting}
\end{frame}

\begin{frame}{Remarks on the code example}
The code is only a template. In a real application you would typically adapt:
\begin{itemize}
\item the recurrent cell (GRU, LSTM, ConvLSTM),
\item the PDE residual,
\item the boundary handling,
\item normalization of the two layers,
\item and the choice of \(\lambda_{\mathrm{phys}}\).
\end{itemize}

For genuine spatiotemporal data, a ConvLSTM or a convolutional recurrent model is often even more natural than flattening the grid.
\end{frame}

\begin{frame}{Final summary}
For a recurrent neural network with two height layers on a fixed \((x,y)\)-grid:
\begin{itemize}
\item use a \textbf{joint multi-output loss} rather than two unrelated losses,
\item weight the two channels if their scales differ,
\item add spatial-gradient penalties if structure matters,
\item and add a physics-informed residual if governing equations are known.
\end{itemize}

A practical default is
\[
\mathcal{L}
=
w_1\mathcal{L}_1 + w_2\mathcal{L}_2 + \lambda_{\mathrm{phys}}\mathcal{L}_{\mathrm{phys}}.
\]
\end{frame}

\end{document}
